\section{Discussion and limitations}

The experiments show that differentiable OED can improve continuous receiver
placement for a finite nonlinear inversion. More broadly, the approach addresses
one source of misspecification in practical FWI: the acquisition geometry must
usually be chosen before the informative measurement locations are known. By
alternating inversion, receiver updates, and reacquisition, the proposed loop
can adapt an initially suboptimal geometry as information about the model
emerges.

The three design objectives should be viewed as different compromises rather than a
definitive ranking. The oracle model error design objective is a privileged reference that is unavailable in a physical experiment. Illumination
currently approaches the oracle’s mean reconstruction benefit at lower computational
cost, but with greater variation across repetitions. Waveform misfit is less
variable but gives a smaller mean improvement. The waveform misfit objective assumes that receiver layouts producing a better data fit after a finite FWI block also produce velocity models closer to the true model. This assumption is not guaranteed in a nonlinear, ill-posed inverse problem.

The present conclusions are limited to noise-free synthetic observations
generated with the same solver used for inversion. Robustness to measurement
noise and modeling error therefore remains untested. Moreover, the methods are
matched by their number of FWI updates rather than by total computational or
physical acquisition cost, and the learned layouts depend on the chosen
optimizers and update budgets. Future work should evaluate independent forward
models, LDV-calibrated noise, fixed acquisition budgets, and ultimately
closed-loop laboratory measurements.

% Illumination avoids differentiation through the unrolled FWI trajectory,
% substantially reducing runtime and peak memory. Its scaling to longer
% trajectories remains untested, and larger problems may require checkpointing,
% truncated unrolling, or implicit gradients.

% The methods are matched by 100 FWI updates rather than total computation or
% physical acquisition time. Future LDV studies should instead compare
% reconstruction quality under fixed acquisition and computational budgets.

% The observations are noise-free and generated with the same solver used for
% inversion, so robustness to noise and modeling error remains untested. Future
% evaluations should use LDV-calibrated noise, an independent forward model, and
% ultimately closed-loop laboratory acquisition.

Overall, the results support differentiable acquisition as a flexible way to
decide not only how to reconstruct a model, but also where the data needed for
that reconstruction should be measured. The present approach provides also an exciting route for
investigating which hyperparameters, optimizers, and wave solvers could
themselves become optimization variables.





% The optimized layouts also depend on the chosen optimizer and update budget;
% these choices could themselves become optimization variables.